\newgeometry{top=3cm, bottom=3cm, left=3cm, right=3cm}
\begin{titlepage}
  \begin{center}
    \vspace*{4cm}
    {\Huge \textbf{Dédicace}}

    \vspace{2cm}

    \begin{minipage}{0.85\textwidth}
      \onehalfspacing

      \textbf{À ma mère.}  
      Ton amour, ta patience et ta confiance m’ont accompagné à chaque étape de ce travail.  
      Tu as toujours été ma source de tendresse et de force.  
      Que Dieu te protège et t’accorde santé, paix et longue vie.

      \vspace{0.3cm}

      \textbf{À mon père.}  
      Ton sens du devoir, ta discrétion et ta présence solide m’ont transmis la rigueur et le courage.  
      Merci pour tous tes sacrifices.  
      Que Dieu te préserve et t’accorde une vie paisible et épanouie.

      \vspace{0.3cm}

      \textbf{À mes frères et à ma sœur.}  
      Vous êtes bien plus qu’une famille : vous êtes mes repères et mes complices.  
      Merci pour votre affection, vos encouragements et votre confiance sans faille.

      \vspace{0.3cm}

      \textbf{À mon oncle Idriss.}  
      Pour ta gentillesse, tes conseils et ta présence. Merci d’avoir cru en moi.

      \vspace{0.3cm}

      \textbf{À ma famille et à mes amis proches.}  
      Merci pour votre soutien constant, vos mots réconfortants et votre bienveillance.  
      Chacun de vous, à sa manière, a contribué à ce projet.

      \vspace{0.8cm}

      \begin{flushright}
        Anas
      \end{flushright}
    \end{minipage}
  \end{center}\end{titlepage}
\restoregeometry

%---------- Remerciements ----------

\begin{figure}[!h]
\begin{center}
\vspace{1.5cm}
\includegraphics[height=1.8cm]{presentation/remerci.png}
\end{center}
\end{figure}
\vspace{0.5cm}

\begin{center}
    \LARGE \textbf{Remerciements}
\end{center}

\vspace{1cm}

\begin{onehalfspacing}

Je tiens d’abord à remercier mes encadrants chez \textbf{N+ONE DATACENTERS}, 
\textbf{M. Ayoub Amellaln} et \textbf{M. Amine El Mahi}, pour leur confiance, 
leur disponibilité et leurs conseils tout au long de ce stage. 
Merci de m’avoir guidé simplement et avec bienveillance, et de m’avoir offert 
un cadre de travail agréable et motivant.

\vspace{0.5cm}

Je remercie aussi le corps enseignant de l’\textbf{ENSIAS} pour la formation 
solide qu’il m’a donnée pendant ces deux ans, 
et tout particulièrement les enseignants de la filière Génie Logiciel.

\vspace{0.5cm}

Enfin, merci à toute l’équipe de \textbf{N+ONE DATACENTERS} pour leur accueil chaleureux 
et pour avoir rendu mon intégration facile dès le premier jour.

\end{onehalfspacing}\vspace{0.3cm}
}
\end{spacing}

\begin{figure}[!h]
\begin{center}
\vspace{1cm}
\includegraphics[height=1cm]{presentation/remerc2.png}
\end{center}
\end{figure}

\newpage

%---------- Résumé ----------

%---------- Résumé ----------

\hskip7mm
\vspace{3cm}
\begin{center}
{\textbf{\Huge Résumé}}
\end{center}
\vspace{1cm}

\begin{spacing}{1.15}

{\large
Ce rapport présente le développement d'une application web sécurisée de gestion
des appels d'offres, réalisé lors d'un stage d'un mois au sein de \textbf{N+ONE DATACENTERS}
à Casablanca. Le projet répond à un besoin concret : remplacer un processus fragmenté
(emails, Excel, partages réseau) par une solution centralisée, traçable et sécurisée.
}
\vspace{0.4cm}

{\large
L'application repose sur une architecture monolithe modulaire (\textbf{Django REST Framework}
côté backend, \textbf{React.js} côté frontend) déployée via \textbf{Docker Compose}
avec un proxy Nginx en HTTPS. Elle couvre l'intégralité du cycle de vie d'un appel
d'offres : de la création au brouillon jusqu'à l'attribution du marché, en passant
par la publication, la collecte des soumissions et l'évaluation des candidats.
}
\vspace{0.4cm}

{\large
Les principales fonctionnalités de sécurité implémentées sont :
l'authentification \textbf{JWT} avec access token en mémoire et refresh token en cookie
\texttt{httpOnly} scopé, le contrôle d'accès basé sur les rôles (\textbf{RBAC}) à quatre
niveaux (ADMIN, MANAGER, ANALYST, VIEWER), le chiffrement sélectif
\textbf{AES-256-GCM} des documents sensibles, et un journal d'audit complet.
Un portail public isolé permet aux sociétés externes de déposer leurs propositions
sans compte utilisateur, avec rate limiting et vérification des deadlines côté serveur.
}
\vspace{0.4cm}

{\large
Un pipeline d'intégration continue (GitHub Actions) exécute automatiquement
les tests pytest et le lint flake8 à chaque push. L'application est déployable
en une commande (\texttt{docker-compose up}).
}
\vspace{0.8cm}

\noindent\rule{\linewidth}{0.4pt}
\vspace{0.3cm}

\noindent
{\large \textbf{Mots-clés :} Appels d'offres, Django REST Framework, React.js,
JWT, RBAC, AES-256-GCM, Docker Compose, GitHub Actions, Portail public, Audit.}

\end{spacing}

\newpage
%---------- Abstract ----------

\hskip7mm
\vspace{3cm}
\begin{center}
{\textbf{\Huge Abstract}}
\end{center}
\vspace{1cm}

\begin{spacing}{1.15}
{\large
This report presents the development of a secure web application for procurement
management, carried out during a one-month internship at \textbf{N+ONE DATACENTERS}
in Casablanca. The project addresses a concrete need: replacing a fragmented process
(emails, Excel, shared drives) with a centralized, traceable, and secure solution.

\vspace{0.4cm}

The application is built on a modular monolith architecture (\textbf{Django REST Framework}
backend, \textbf{React.js} frontend) deployed via \textbf{Docker Compose} with an Nginx
HTTPS proxy. It covers the complete procurement lifecycle: from draft creation to contract
award, through publication, submission collection, and candidate evaluation.

\vspace{0.4cm}

The main security features implemented are: \textbf{JWT authentication} with access token
in memory and refresh token in a scoped \texttt{httpOnly} cookie, four-level
\textbf{role-based access control} (RBAC), selective \textbf{AES-256-GCM encryption}
of sensitive documents, and a comprehensive audit log. An isolated public portal allows
external companies to submit proposals without an account, with IP-based rate limiting
and server-side deadline enforcement.

\vspace{0.4cm}

A continuous integration pipeline (GitHub Actions) automatically runs pytest tests and
flake8 lint on every push. The application deploys in a single command
(\texttt{docker-compose up}).

\vspace{0.8cm}

\noindent\rule{\linewidth}{0.4pt}
\vspace{0.3cm}

\noindent\textbf{Keywords:} Procurement management, Django REST Framework, React.js,
JWT, RBAC, AES-256-GCM, Docker Compose, GitHub Actions, Public portal, Audit log.
}
\end{spacing}

\newpage

%---------- Liste des abréviations ----------

\chapter*{Liste des abréviations}
\addcontentsline{toc}{chapter}{Liste des abréviations}

\begin{acronym}[AES-256-GCM]
  \acro{AES}{Advanced Encryption Standard}
  \acro{API}{Application Programming Interface}
  \acro{CI}{Continuous Integration}
  \acro{CORS}{Cross-Origin Resource Sharing}
  \acro{CRUD}{Create, Read, Update, Delete}
  \acro{DRF}{Django REST Framework}
  \acro{GCM}{Galois/Counter Mode}
  \acro{HTTP}{HyperText Transfer Protocol}
  \acro{HTTPS}{HyperText Transfer Protocol Secure}
  \acro{JWT}{JSON Web Token}
  \acro{MIME}{Multipurpose Internet Mail Extensions}
  \acro{ORM}{Object-Relational Mapper}
  \acro{PDF}{Portable Document Format}
  \acro{RBAC}{Role-Based Access Control}
  \acro{REST}{Representational State Transfer}
  \acro{SPA}{Single Page Application}
  \acro{SQL}{Structured Query Language}
  \acro{TLS}{Transport Layer Security}
  \acro{UUID}{Universally Unique Identifier}
  \acro{XSS}{Cross-Site Scripting}
\end{acronym}

\newpage
